\section{Decidable and computably enumerable sets}
\label{sec:turing-machines-03}

\begin{definition}[Decidable set]
A set \(A\subseteq\N\) is \emph{decidable} iff its characteristic function
\(\chi_A:\N\to\{0,1\}\) is computable.
\end{definition}

\begin{definition}[Computably enumerable set]
A set \(A\subseteq\N\) is \emph{computably enumerable} iff it is the domain
of a partial computable function.
\end{definition}

\begin{definition}[Decidable and computably enumerable relations]
For \(R\subseteq\N^k\), the relation \(R\) is \emph{decidable} iff its
characteristic function is computable.  It is \emph{computably enumerable}
iff
\[
  \set{\operatorname{code}_k(\bar n):\bar n\in R}
\]
is a computably enumerable subset of \(\N\).  These terms classify relations;
for \(k=1\) they reproduce the preceding definitions for sets.  The extension
to relations is needed by the computation-history and proof predicates below.
\end{definition}

\begin{definition}[Machine decision and recognition]
Let \(A\subseteq\N\).  A Turing machine \emph{decides} \(A\) iff it halts on
every unary numerical input \(\operatorname{bin}(n)\), accepts exactly those inputs
whose decoded number \(n\) belongs to \(A\), and rejects exactly those whose
decoded number lies outside \(A\).  A Turing machine \emph{recognizes} \(A\)
iff it accepts exactly the unary numerical inputs whose decoded numbers belong
to \(A\), with no halting requirement on inputs outside \(A\).  A machine with
the first property is a \emph{decider} for \(A\); a machine with the second is
a \emph{recognizer} for \(A\).

These terms classify machine behavior, not sets or numerical functions.  They
are required to state the machine forms of decidability and computable
enumerability without ambiguity.  Every decider is a recognizer, but a
recognizer may diverge outside \(A\), so the two terms are not redundant.  The
next lemma and theorem are their uses.
\end{definition}

\begin{lemma}[Machine form of decidability]
A set \(A\subseteq\N\) is decidable iff some Turing machine decides \(A\).
\textit{Status: complete proof.}
\end{lemma}

\begin{proof}
Given a machine computing \(\chi_A\), simulate it and accept when its output is
\(1\), while rejecting when its output is \(0\).  Totality makes this machine
halt on every input.  Conversely, simulate a decider and write \(1\) after an
accepting halt and \(0\) after a rejecting halt.  The resulting machine
computes \(\chi_A\).
\end{proof}

\begin{theorem}[Equivalent forms of computable enumerability]
For \(A\subseteq\N\), the following are equivalent:
\begin{enumerate}
  \item \(A\) is computably enumerable;
  \item some Turing machine recognizes \(A\);
  \item either \(A=\varnothing\), or \(A\) is the range of a computable
        function \(e:\N\to\N\); and
  \item there is a decidable relation \(R(x,k)\) such that
        \(x\in A\) iff \(\exists k\,R(x,k)\).
\end{enumerate}
\textit{Status: proof sketch.}
\end{theorem}

\begin{proofplan}
Convert a domain machine into a recognizer.  Dovetail all computations to
enumerate the accepted inputs.  Record a finite accepting computation as the
witness \(k\) for a decidable computation-history relation.  Conversely,
search successive \(k\) for a witness to \(R(x,k)\).  The empty-set exception
in the range formulation is necessary because a total function on \(\N\) has
nonempty range.
\end{proofplan}

\begin{theorem}
A set \(A\subseteq\N\) is decidable iff both \(A\) and
\(\N\setminus A\) are computably enumerable.
\textit{Status: complete proof.}
\end{theorem}

\begin{proof}
If \(A\) is decidable, a decider recognizes \(A\), and reversing its
two halting answers gives a recognizer for the complement.  Conversely, run a
recognizer for \(A\) and a recognizer for its complement in dovetailing fashion.
Exactly one eventually accepts, and its identity determines membership in
\(A\).
\end{proof}

